% !PRINTED 251
\weberpage{251}
%<<CONT
and from this there follows, having regard to
\[
\gamma_3(\omega+1) = -\gamma_3(\omega), \quad
\gamma_3\left(-\frac{1}{\omega}\right) = -\gamma_3(\omega): \qquad \text{[§\,54, (14)]}
\]

3. \gesperrt{If $n$ is odd and $c$ divisible by 4, then the
quantities}
\[
P^{6}_{c,\eth,a}\,\gamma_3(\omega)^{\frac{n-1}{2}}
\]
\gesperrt{are roots of a multiplier equation.} If $n \equiv 1 \pmod{4}$, the same holds of $P^{6}_{c,\eth,a}$.

As an example we choose $n=3$ and set

\weq{5}{x = 27\left(\frac{\eta(3\omega)}{\eta(\omega)}\right)^6, \quad
x_0 = -\left(\frac{\eta\left(\dfrac{\omega}{3}\right)}{\eta(\omega)}\right)^6, \quad
x_1 = -\left(\frac{\eta\left(\dfrac{4+\omega}{3}\right)}{\eta(\omega)}\right)^6, \quad
x_2 = -\left(\frac{\eta\left(\dfrac{8+\omega}{3}\right)}{\eta(\omega)}\right)^6.}

The coefficients of the equation which has these quantities for its
roots are rational functions of $\gamma_3(\omega)$, and since none of
the quantities (5) becomes infinite for a finite value of $\gamma_3$,
they are \gesperrt{polynomials} in $\gamma_3$. According to 3. we may
therefore set this equation up in the form:
\[
x^4 + A_1\gamma_3 x^3 + A_2 x^2 + A_3 \gamma_3 x + A_4 = 0,
\]
in which $A_1, A_2, A_3, A_4$ are polynomials in
$j(\omega)$. First of all one obtains $A_4$ as the product $x x_0 x_1 x_2$,
which can vanish for no value of $j(\omega)$ and must therefore be
constant. From the beginnings of the development [§\,24]

\weq{6}{\begin{aligned}
x &= 27\,q\dots,\\
x_0 &= -q^{-\frac{1}{3}}\dots,\\
x_1 &= -e^{\frac{2\pi i}{3}}q^{-\frac{1}{3}}\dots,\\
x_2 &= -e^{-\frac{2\pi i}{3}}q^{-\frac{1}{3}}\dots,\\
\gamma_3 &= q^{-1}\dots
\end{aligned}}

one therefore finds
\[
A_4 = -27.
\]
Further we have
\[
-A_1\gamma_3 = x + x_0 + x_1 + x_2.
\]
Since the right-hand side does not become infinite even in the first
order for an infinite $\gamma_3$, that is, for $q=0$, $A_1$ must
vanish.
